\chapter{On-disk format}
\label{ch:ondisk}

\section{0: Purpose}
This document defines the byte-level contract shared by mkfs, bitter-dump, bitter-fsck, an the kernel module.
\lstinline|format/bitterfs_format.h| is authorative and will be explained in this document. 
\section{1: Disk Format and Conventions}
\bf{This section outlines the formatting and conventions that will be used in this file system}
Byte order is little endian, with \lstinline|__le32/__le64| on disk. 
Accessors will be used to convert that format to the 
host CPU's byte order.
Block size is 4096 a common standard across FS design
Addressing happens with byte offsets("bytenr") which are always a multiple of 4096
checksum 32 byte field; crc32c in the first 4 bytes, rest zero
checksum covers every after the checksum field to the end of the block
\section{2: Layout map}
\bf{Outlines the byte add}
0x00000 is reserved 64 KiB, that must remain untouched
0x10000 is a superblock, it is the only fixed block in the entire file system
0x11000 is the initial address of the rootnode
0x12000 is the first free block set at initialization
Everything past this point is free
\section{3: Objectid map}


Object ids are 64-bit and mostly assigned freely to files and directories, but
both ends of the range are reserved. Reservations cannot be added later without
renumbering every image ever made, which is why they are settled before
\texttt{mkfs} exists.

\begin{center}
\begin{tabular}{llp{7.6cm}}
  \toprule
  Value & Name & Meaning \\
  \midrule
  0     & \texttt{INVALID}     & Never assigned, so that an all-zero key is recognisably invalid \\
  1     & \texttt{ROOT\_TREE}   & The tree of tree roots, located via \texttt{super.root} \\
  2     & \texttt{EXTENT\_TREE} & Allocation records and reference counts \\
  4     & \texttt{FS\_TREE}     & Inodes, directory entries and extent maps \\
  256   & \texttt{FIRST\_FREE}  & First id available to files and directories \\
  \bottomrule
\end{tabular}
\end{center}

\decide{The top of the object id space is not yet reserved. btrfs keeps the
region at \texttt{-256} and above for special objects. Reserving it costs
nothing now and is impossible later.}

\section{Item types}
\label{sec:item-types}

The \texttt{type} byte does two jobs: it says how to interpret the payload, and
it says what the key's \texttt{offset} field means.

\begin{center}
\begin{tabular}{llll}
  \toprule
  Value & Name & \texttt{objectid} holds & \texttt{offset} holds \\
  \midrule
  1   & \texttt{INODE\_ITEM}   & inode number  & 0, unused \\
  25  & \texttt{INODE\_REF}    & inode number  & the parent's inode number \\
  48  & \texttt{DIR\_ITEM}     & parent inode  & hash of the entry's name \\
  72  & \texttt{DIR\_INDEX}    & parent inode  & per-directory sequence number \\
  96  & \texttt{EXTENT\_DATA}  & inode number  & offset within the file \\
  120 & \texttt{EXTENT\_ITEM}  & disk address  & length of the extent \\
  255 & \texttt{KEY\_TYPE\_MAX} & \multicolumn{2}{l}{reserved sentinel, never a real type} \\
  \bottomrule
\end{tabular}
\end{center}

\subsection{Why these numbers}

The values are not labels. Keys sort by \texttt{objectid}, then \texttt{type},
then \texttt{offset}, so \textbf{the numeric order of these constants is the
physical order items appear on disk} for a given object. Only the relative
order carries design information: renumbering everything consistently would
change nothing, but swapping two values changes the layout.

\texttt{INODE\_ITEM} is 1 so an object's inode is always the first item in its
range, which any lookup reaches immediately. \texttt{EXTENT\_DATA} is last of
the file types because a large file has thousands of them, and that bulk
belongs at the end of the range rather than between the inode and its directory
entries. It also makes \texttt{truncate} a single contiguous range delete that
runs to the end of the object's keys with nothing after it to preserve.

The gaps are insertion room. A type added later --- \texttt{EXTENT\_CSUM} for
data checksums, or \texttt{XATTR\_ITEM} --- must be able to land at the correct
\emph{sort position} without renumbering anything already on disk.

\subsection{The sentinel}

255 is reserved because range scans need an upper bound that sorts after every
genuine item. To visit everything belonging to object $N$, scan from
\texttt{(N, 0, 0)} to \texttt{(N, 255, (bt\_le64)-1)}. Were 255 assignable, the
last item of that type would fall outside the range and \texttt{unlink} would
leak it.

\section{Structures}
\label{sec:structures}

Offsets below are read from the compiled header rather than computed by hand,
so this table and the code cannot disagree. Every structure is packed: padding
is explicit and named, never compiler-inserted.

\subsection{\texttt{bitter\_key}}

\begin{center}
\begin{bytefield}[bitwidth=3.1em]{8}
  \bitheader{0-7} \\
  \wordbox{1}{objectid \hfill \texttt{bt\_le64}} \\
  \bitbox{1}{type} & \bitbox{7}{offset (bytes 0--6)} \\
  \bitbox{1}{\dots} & \bitbox{7}{\emph{next structure begins here}} \\
\end{bytefield}
\end{center}

\begin{center}
\begin{tabular}{rrll}
  \toprule
  Offset & Size & Type & Field \\
  \midrule
  0 & 8 & \texttt{bt\_le64} & \texttt{objectid} \\
  8 & 1 & \texttt{bt\_u8}   & \texttt{type} \\
  9 & 8 & \texttt{bt\_le64} & \texttt{offset} \\
  \midrule
  \multicolumn{2}{r}{\textbf{17}} & & \\
  \bottomrule
\end{tabular}
\end{center}

Seventeen bytes, deliberately not a multiple of anything. The diagram above
makes the consequence visible: the second key in an array starts one byte into
a row. That is the price of packing, and it is worth paying --- the layout is
then identical on every target, with no compiler free to insert padding of its
own choosing.

Note that \texttt{type} is \texttt{bt\_u8} rather than a little-endian type. A
single byte has no byte order, so it needs no conversion and no marking.

\subsection{\texttt{bitter\_header}}

Begins every tree block, internal node and leaf alike.

\begin{center}
\begin{tabular}{rrllp{6.2cm}}
  \toprule
  Offset & Size & Type & Field & Purpose \\
  \midrule
  0  & 32 & \texttt{bt\_u8[]}  & \texttt{csum}       & crc32c in bytes 0--3, rest zero; covers everything after this field \\
  32 & 16 & \texttt{bt\_u8[]}  & \texttt{fsid}       & must equal the superblock's \\
  48 & 8  & \texttt{bt\_le64}  & \texttt{bytenr}     & this block's own address \\
  56 & 8  & \texttt{bt\_le64}  & \texttt{generation} & the transaction that wrote it \\
  64 & 8  & \texttt{bt\_le64}  & \texttt{owner}      & which tree this block belongs to \\
  72 & 4  & \texttt{bt\_le32}  & \texttt{nritems}    & items in a leaf, children in a node \\
  76 & 1  & \texttt{bt\_u8}    & \texttt{level}      & 0 means leaf \\
  77 & 19 & \texttt{bt\_u8[]}  & \texttt{reserved}   & zero on write \\
  \midrule
  \multicolumn{2}{r}{\textbf{96}} & & & \\
  \bottomrule
\end{tabular}
\end{center}

\paragraph{Three fields make a block self-identifying.} \texttt{csum} proves the
contents are intact; \texttt{fsid} proves the block belongs to \emph{this}
filesystem; \texttt{bytenr} proves it was found where it claims to live. The
last two catch what a checksum structurally cannot --- a stale block from a
previous \texttt{mkfs}, or a misdirected write. Both checksum perfectly,
because both were once valid blocks of something else.

\texttt{owner} closes the remaining gap: a block that is intact, in the right
place, and belongs to this filesystem, but is claimed by the wrong tree.

\paragraph{\texttt{generation} is the copy-on-write witness.} A block whose
generation equals the current transaction is already ours and may be modified
in place; anything older must be copied first. That single comparison is the
entire CoW decision.

\paragraph{Why 96 and not 77.} A multiple of 16, so structures align with
hexdump rows, and it leaves the leaf data area at exactly 4000 bytes --- which
makes capacity arithmetic something that can be done in the head while
debugging. btrfs by contrast uses an unpadded 101-byte header, and every
hexdump of a btrfs leaf is correspondingly awkward to read.

\subsection{\texttt{bitter\_key\_ptr}}

An entry in an internal node.

\begin{center}
\begin{tabular}{rrllp{6cm}}
  \toprule
  Offset & Size & Type & Field & Purpose \\
  \midrule
  0  & 17 & \texttt{bitter\_key} & \texttt{key}        & smallest key in the subtree below \\
  17 & 8  & \texttt{bt\_le64}    & \texttt{blockptr}   & address of the child block \\
  25 & 8  & \texttt{bt\_le64}    & \texttt{generation} & child's generation, cross-checked on read \\
  \midrule
  \multicolumn{2}{r}{\textbf{33}} & & & \\
  \bottomrule
\end{tabular}
\end{center}

At 33 bytes per entry a node holds $\lfloor 4000/33 \rfloor = 121$ children, so
three levels address roughly $121^3 \approx 1.7$ million leaves. Tree depth
grows slowly enough that 4096-byte nodes remain comfortable at this scale.

\subsection{\texttt{bitter\_item}}

An entry in a leaf.

\begin{center}
\begin{tabular}{rrllp{6cm}}
  \toprule
  Offset & Size & Type & Field & Purpose \\
  \midrule
  0  & 17 & \texttt{bitter\_key} & \texttt{key}    & \\
  17 & 4  & \texttt{bt\_le32}    & \texttt{offset} & payload start, measured from the end of the header \\
  21 & 4  & \texttt{bt\_le32}    & \texttt{size}   & payload length \\
  \midrule
  \multicolumn{2}{r}{\textbf{25}} & & & \\
  \bottomrule
\end{tabular}
\end{center}

\paragraph{Leaves grow from both ends.} These descriptors march forward from the
header while the payloads they point at stack backward from the end of the
block; the free space is the gap between. That is why \texttt{offset} is
measured from the end of the header --- state the base explicitly, because both
conventions exist and mixing them is a guaranteed bug.

\decide{\texttt{offset} and \texttt{size} are 32 bits each, inherited from
btrfs, which supports node sizes up to 64\,KiB. Both values are bounded by a
4096-byte block here, so they need only 12 bits. \texttt{uint16\_t} would make
the item 21 bytes instead of 25 --- about 19\% more items per leaf, and a
correspondingly shallower tree. The argument for keeping 32 bits is that every
intermediate in the leaf arithmetic is then obviously overflow-free, and that
arithmetic is where the bugs live.}

\subsection{\texttt{bitter\_root\_item}}

Where a tree's root block currently lives. The payload of a
\itemtype{ROOT\_ITEM}, stored under \keytriple{objectid}{ROOT\_ITEM}{0} in the
root tree, one per tree.

\begin{center}
\begin{tabular}{rrllp{6cm}}
  \toprule
  Offset & Size & Type & Field & Purpose \\
  \midrule
  0  & 8 & \texttt{bt\_le64} & \texttt{bytenr}     & the tree's root block \\
  8  & 8 & \texttt{bt\_le64} & \texttt{generation} & that block's generation \\
  16 & 1 & \texttt{bt\_u8}   & \texttt{level}      & the tree's height \\
  \midrule
  \multicolumn{2}{r}{\textbf{17}} & & & \\
  \bottomrule
\end{tabular}
\end{center}

\paragraph{No \field{objectid}.} The key already carries it. Two copies could
disagree, and the one nobody reads is the one that rots.

\paragraph{The root tree has no root item.} It cannot: the item would have to
live inside the tree it locates. Its address comes from \field{super.root} by
fiat, and every other tree is found through it. That asymmetry is the only
place in the format where a lookup bottoms out rather than recursing.

\paragraph{\field{level} is duplicated} --- it is also in the root block's own
header. A descent needs the depth \emph{before} the first read, which is the
same argument that put \field{root\_level} in the superblock.

\decide{\field{offset} in the key is 0 and means ``unused''. btrfs stores the
creating transaction there, so a tree's roots sort in creation order and
``find the newest'' is a backward scan from \code{(u64)-1}. Nothing needs that
until snapshots exist. Whatever the field comes to mean, 0 must keep meaning
the live root, because every image already written has a 0 in it.}

\subsection{\texttt{bitter\_extent\_item}}

How many references point at one extent. Stored under
\keytriple{bytenr}{EXTENT\_ITEM}{length}.

\begin{center}
\begin{tabular}{rrllp{6cm}}
  \toprule
  Offset & Size & Type & Field & Purpose \\
  \midrule
  0 & 8 & \texttt{bt\_le64} & \texttt{refs}  & pointers at this extent \\
  8 & 8 & \texttt{bt\_le64} & \texttt{flags} & what kind of extent it is \\
  \midrule
  \multicolumn{2}{r}{\textbf{16}} & & & \\
  \bottomrule
\end{tabular}
\end{center}

\begin{center}
\begin{tabular}{llp{7.5cm}}
  \toprule
  Bit & Name & Meaning \\
  \midrule
  0 & \code{BITTER\_EXTENT\_FLAG\_RESERVED} &
      allocated and reachable from no tree: the reserved prefix and the
      superblock \\
  \bottomrule
\end{tabular}
\end{center}

\paragraph{The key carries the identity.} The extent's start address is the
objectid and its length is the offset, so the payload holds only what changes.
One counted field is the whole structure for that reason, not for want of
finishing it.

\paragraph{An extent is free when it has \emph{no} item}, not when it has an
item holding zero. The last reference deletes the item, so a stored 0 is
corruption. This matters because free space is derived by subtraction: the gaps
between items are what is available, and an item that lingered at zero would
make free space look allocated forever.

\paragraph{Address order is not an accident.} Keys sort
$(\mbox{objectid},\,\mbox{type},\,\mbox{offset})$ and the objectid here \emph{is}
an address, so every extent item in the filesystem sits in the tree in device
order. Rebuilding the free-space map is therefore one linear walk with no
sorting step --- see \cref{ch:extents}.

\paragraph{Why \code{RESERVED} exists.} Without it, the 68\,KiB before the
first tree block is an extent with one reference that nothing points at, which
is indistinguishable from a leak. \fnname{bitter-fsck} would report it on every
filesystem that has ever been created, and no repair could ever clear it.

\decide{An unknown flag bit must be preserved rather than cleared, so a newer
version's extents survive an older tool. Deliberately absent: \field{generation}
(nothing reads it), and the tree-block-versus-file-data distinction, which waits
for phase 7 to have data extents to distinguish.}

\subsection{\texttt{bitter\_super}}

\begin{center}
\begin{tabular}{rrllp{5.4cm}}
  \toprule
  Offset & Size & Type & Field & Purpose \\
  \midrule
  0   & 32  & \texttt{bt\_u8[]} & \texttt{csum}            & covers everything after this field \\
  32  & 16  & \texttt{bt\_u8[]} & \texttt{fsid}            & the filesystem's UUID \\
  48  & 8   & \texttt{bt\_u8[]} & \texttt{magic}           & identifies the image as bitterfs \\
  56  & 8   & \texttt{bt\_le64} & \texttt{bytenr}          & its own address, \texttt{0x10000} \\
  64  & 8   & \texttt{bt\_le64} & \texttt{generation}      & the committed transaction number \\
  72  & 8   & \texttt{bt\_le64} & \texttt{root}            & address of the root tree's root block \\
  80  & 8   & \texttt{bt\_le64} & \texttt{total\_bytes}    & device size recorded at mkfs \\
  88  & 8   & \texttt{bt\_le64} & \texttt{bytes\_used}     & allocated bytes, for \texttt{statfs} \\
  96  & 8   & \texttt{bt\_le64} & \texttt{compat\_flags}   & unknown bit: mount rw normally \\
  104 & 8   & \texttt{bt\_le64} & \texttt{ro\_compat\_flags} & unknown bit: mount read-only \\
  112 & 8   & \texttt{bt\_le64} & \texttt{incompat\_flags} & unknown bit: refuse to mount \\
  120 & 4   & \texttt{bt\_le32} & \texttt{block\_size}     & 4096 \\
  124 & 2   & \texttt{bt\_le16} & \texttt{csum\_type}      & which algorithm fills \texttt{csum} \\
  126 & 1   & \texttt{bt\_u8}   & \texttt{root\_level}     & depth of the root tree \\
  127 & 1   & \texttt{bt\_u8}   & \texttt{reserved0}       & alignment, so \texttt{label} starts at 128 \\
  128 & 32  & \texttt{bt\_u8[]} & \texttt{label}           & human-readable name for \texttt{blkid} \\
  160 & 352 & \texttt{bt\_u8[]} & \texttt{reserved}        & future room \\
  \midrule
  \multicolumn{2}{r}{\textbf{512}} & & & \\
  \bottomrule
\end{tabular}
\end{center}

\paragraph{The only structure never copy-on-written.} There is nowhere to put a
pointer to a new copy of the superblock, so it is overwritten in place every
commit. That single overwrite \emph{is} the transaction commit: everything else
is copy-on-written precisely so that this one write can be the moment the
transaction becomes real.

\paragraph{Why exactly 512 bytes.} If the structure spanned two sectors, a torn
write could leave the first new and the second old. The checksum would then
fail, and with a single superblock copy the filesystem would be unmountable.
Padding to exactly one sector makes atomicity structural rather than hoped for,
and \texttt{\_Static\_assert} enforces it at compile time.

\paragraph{Field order.} Largest first, which incidentally lands every 64-bit
field on an 8-byte boundary, \texttt{block\_size} on a 4-byte one and
\texttt{label} on 16. Packing removes the compiler's right to insert padding; it
does not prevent laying fields out sensibly by hand.

\paragraph{Feature flags rather than a version number.} A single version integer
forces a total ordering on features: version 3 must contain everything from 1
and 2. Three flag words allow independent features to be added independently,
and the read-only case is the interesting one. Consider adding data checksums
later: an old reader can still read every file correctly, but were it to
\emph{write} one it would not create the checksum items, and a newer reader
would see data with missing checksums and call it corruption. Read-only is
exactly the right restriction, and no version number can express it.

\section{Derived constants}
\label{sec:derived}

Computed from the structure sizes rather than written down, so that changing
the header padding updates them automatically and a stale figure cannot survive
in the code.

\begin{center}
\begin{tabular}{lrl}
  \toprule
  Constant & Value & Definition \\
  \midrule
  \texttt{BITTER\_LEAF\_DATA\_SIZE} & 4000 & block size $-$ \texttt{sizeof(bitter\_header)} \\
  items per leaf                    & 160  & $4000 / 25$ \\
  children per node                 & 121  & $4000 / 33$ \\
  \bottomrule
\end{tabular}
\end{center}
